\begin{center}
    \section*{\fontsize{20}{20}\selectfont Chapter 2}
\end{center}
\vspace{10mm}
\section{Background Study and Literature Review}

Brain tumor segmentation has come a long way over the past two decades. Advances in medical imaging and deep learning have pushed the field forward significantly. Accurately outlining tumor boundaries from multi-parametric MRI scans matters for diagnosis, treatment planning, and tracking patients over time. This chapter walks through how automated segmentation methods have evolved—from classical image processing to the CNN revolution, and now toward graph-based approaches. Along the way, we'll see why computational efficiency and clinical deployment remain persistent challenges, which is exactly what motivates the GNN framework proposed in this thesis.

\subsection{Evolution of Medical Image Segmentation}

\subsubsection{Classical Approaches}

Early brain tumor segmentation depended on traditional computer vision techniques: intensity thresholding, region growing, atlas-based methods. These required lots of manual tuning and struggled with how differently tumors can look from patient to patient. The BraTS Challenge, which started in 2012 \cite{menze2015multimodal}, changed things by standardizing evaluation protocols and providing large-scale annotated datasets. That really accelerated innovation in the field.

\subsubsection{The CNN Revolution: U-Net and Its Variants}

Fully convolutional networks really transformed medical image segmentation. When Ronneberger et al. \cite{ronneberger2015unet} introduced U-Net in 2015, it was a game-changer. The symmetric encoder-decoder architecture with skip connections preserved spatial information during upsampling, which proved crucial. U-Net became the go-to standard for medical segmentation, achieving way better performance than classical methods while needing minimal training data—thanks to aggressive data augmentation.

Taking U-Net to 3D medical imaging, \c{C}i\c{c}ek et al. \cite{cicek20163d} extended it with 3D convolutions, which enabled direct processing of multi-slice MRI scans. But there's a catch: 3D U-Net's computational requirements scale cubically with volume dimensions. That means you need substantial GPU memory—8GB+ VRAM—and that limits deployment on devices with fewer resources.

Then came nnU-Net from Isensee et al. \cite{isensee2021nnunet}, a self-configuring framework that automatically optimizes preprocessing, architecture, and training hyperparameters for whatever medical segmentation task you throw at it. It achieves state-of-the-art results across diverse datasets. However, the automatic configuration needs extensive computational resources for hyperparameter search, and the final models stay large (80M+ parameters).

\subsubsection{Transformer-Based Architectures}

Recent works have explored vision transformers for medical segmentation. Hatamizadeh et al. \cite{hatamizadeh2022swin} proposed Swin-UNETR, combining Swin Transformer encoders with CNN decoders, achieving 93.8\% Dice on BraTS 2021. TransBTS \cite{wang2021transbts} similarly integrates transformer self-attention with 3D convolutions. While these hybrid architectures push accuracy boundaries, they require even larger computational budgets (92M+ parameters, 24GB+ VRAM), exacerbating deployment challenges for clinical settings.

\subsection{Graph Neural Networks in Medical Imaging}

\subsubsection{Fundamentals of Graph Representation Learning}

Graph neural networks extend deep learning to non-Euclidean data structures, learning representations through iterative message passing between connected nodes \cite{hamilton2017graphsage, kipf2017gcn}. Hamilton et al. \cite{hamilton2017graphsage} introduced GraphSAGE (Graph Sample and Aggregate), which learns node embeddings by sampling and aggregating features from local neighborhoods. The inductive learning capability of GraphSAGE enables generalization to unseen graph structures, making it suitable for medical imaging where each patient represents a distinct graph.

\subsubsection{Superpixel Segmentation}

Superpixel algorithms group perceptually similar pixels into compact regions, providing mid-level representations that preserve object boundaries while reducing dimensionality. Achanta et al. \cite{achanta2012slic} developed Simple Linear Iterative Clustering (SLIC), a computationally efficient method that performs k-means clustering in combined color-spatial space. SLIC generates uniform superpixels with low computational cost (linear in pixel count) and controllable compactness, making it ideal for large-scale medical image preprocessing.

\subsubsection{Graph-Based Medical Segmentation: Existing Work}

Several studies have explored graph representations for medical imaging:

\textbf{Hybrid CNN-GNN Approaches:} Most existing work combines CNNs for feature extraction with GNNs for spatial reasoning. These hybrid methods retain the computational overhead of volumetric convolutions while adding graph processing costs, yielding limited efficiency gains.

\textbf{Region Adjacency Graphs:} Some methods construct graphs where nodes represent anatomical structures (organs, lobes) and edges encode spatial relationships. While interpretable, these approaches require prior segmentation of regions, introducing preprocessing dependencies.

\textbf{Point Cloud Representations:} Voxel-to-point sampling combined with PointNet-style architectures has been investigated, but point clouds discard spatial structure information that graphs can preserve through explicit edge connectivity.

\subsection{Research Gap and Motivation}

\subsubsection{The Efficiency-Accuracy Dilemma}

Despite impressive accuracy numbers, modern deep learning methods hit persistent deployment barriers:

\textbf{Computational Requirements:} State-of-the-art models like nnU-Net, Swin-UNETR, and TransBTS need 68M-92M parameters and 8GB-24GB of GPU memory. That limits deployment to high-end research hardware. Rural clinics, mobile diagnostic units, and developing regions just don't have access to that kind of infrastructure.

\textbf{Inference Latency:} Processing a full 240$\times$240$\times$155 MRI volume through deep 3D CNNs takes 50-100ms per patient. That's a problem for real-time clinical workflows where radiologists expect immediate feedback.

\textbf{Energy Consumption:} Large models burn through power during inference. As medical AI scales to millions of patients globally, this raises environmental concerns. Sustainable healthcare AI needs architectures that balance accuracy with resource efficiency.

\subsubsection{Unexplored GNN Potential}

Graph neural networks have proven themselves in social networks, molecular modeling, and recommendation systems. But their application to medical image segmentation? Still pretty new:

\textbf{Pure GNN Approaches Lacking:} Most medical imaging GNN research uses hybrid CNN-GNN designs. \textit{Nobody's rigorously benchmarked a pure, lightweight GNN architecture for efficient brain tumor segmentation.}

\textbf{Superpixel-GNN Integration Underexplored:} Combining superpixel preprocessing (890$\times$ dimensionality reduction) with inductive graph learning (GraphSAGE) is an unexplored avenue. It could potentially achieve competitive accuracy with dramatic efficiency improvements.

\textbf{Reproducibility Gap:} Medical AI research suffers from reproducibility issues—insufficient documentation, unreported hyperparameters, potential data leakage. Establishing rigorous validation standards (patient-level splits, comprehensive integrity checks) is critical for the field to mature.

\subsection{Positioning This Work}

This thesis addresses the identified research gap through three key contributions:

\textbf{1. Pure GNN Architecture:} We demonstrate that a lightweight GraphSAGE model (439K parameters) without any convolutional components can achieve 92.92\% Dice—approaching state-of-the-art transformers while offering 156$\times$ parameter reduction and 6.9$\times$ inference speedup.

\textbf{2. Rigorous Validation Framework:} Patient-level stratified cross-validation with 15 comprehensive integrity checks ensures zero data leakage and establishes reproducibility standards for the field.

\textbf{3. Clinical Deployment Focus:} By achieving real-time performance (12.7ms per patient) on consumer-grade hardware (RTX 2060 6GB), we demonstrate practical viability for resource-constrained healthcare settings.

\vspace{5mm}

This literature review establishes that while CNNs and transformers have pushed accuracy boundaries for brain tumor segmentation, their computational demands create deployment barriers that limit clinical impact. Graph neural networks, through exploitation of tumor sparsity via superpixel representations, offer a promising yet underexplored paradigm for efficient medical image analysis. The following chapters detail our systematic investigation of this approach, demonstrating that competitive accuracy need not require massive models when architectural design aligns with domain-specific structural properties.

\clearpage